\subsection{HAUSDORFF QUOTIENT SPACES}

Let us look for conditions under which a quotient space X/R is Hausdorff (in which case the equivalence relation R is said to be Hausdorff). In the first place, if X/R is Hausdorff, then the subsets of X/R consisting of a single point are closed (no. 1, Proposition 4) and hence each equivalence class mod R is closed in X. But this necessary condition is not sufficient.

The definition of open sets in X/R gives rise to the following necessary and sufficient condition: X/R is Hausdorff if and only if any two distinct equivalence classes in X are contained in disjoint saturated open subsets of X. We shall give other more usable conditions.

PROPOSITION 8. A necessary condition for the quotient space X/R to be Hausdorff is that the graph C of R is closed in X × X. If the equivalence relation R is open, this condition is also sufficient.

Let \(\varphi: X \to X/R\) be the canonical mapping; then \(C\) is the inverse image under \(\varphi \times \varphi: X \times X \to (X/R) \times (X/R)\) of the diagonal \(\Delta\) of \((X/R) \times (X/R)\) . The first part of the proposition therefore follows from the continuity of \(\varphi \times \varphi\) {[}Axiom \((H^{ii})\) and Theorem 1 of §2, no. 1{]}. If \(R\) is open then \((X/R) \times (X/R)\) can be identified with the quotient space \((X \times X)/(R \times R)\) (§5, no. 3, corollary to Proposition 8), and \(\Delta\) is then identified with the canonical image in \((X \times X)/(R \times R)\) of the set \(C\) , which is saturated with respect to \(R \times R\) . Hence \(\Delta\) is closed in \(X \times X\) and therefore \(X\) is Hausdorff.

\pandocbounded{\includegraphics[keepaspectratio]{images/part001_7b053cb2.jpg}}

If R is not open, there are examples where C is closed but R is not Hausdorff (Exercises 10 and 28).

To show that X/R is Hausdorff we can also apply Proposition 5 of no. 1: M and N being two distinct equivalence classes of R it is sufficient that there should be a continuous mapping f of an open subset A of X, saturated with respect to R and containing M and N, into a Hausdorff space \(X'\) , such that 1) f is constant on each equivalence class mod R contained in A, 2) f takes distinct values on M and N. For since \(A/R_{A}\) can be identified with an open subset of X/R (§ 3, no. 6, Proposition 10, Corollary 1), we can apply Proposition 5 of no. 1 to the mapping g: \(A/R_{A} \to X'\) induced by f, since g is continuous (§ 3, no. 4, Proposition 6).

In particular:

PROPOSITION 9. If \(f\) is a continuous mapping of a topological space \(\mathbf{X}\) into a Hausdorff space \(\mathbf{X}'\) , and \(\mathbf{R}\) is the equivalence relation \(f(x) = f(y)\) , then the quotient space \(\mathbf{X} / \mathbf{R}\) is Hausdorff.

PROPOSITION 10. If X is a Hausdorff space, and if X has a continuous section s with respect to the equivalence relation R, then X/R is Hausdorff and s (X/R) is closed in X.

For (§ 3, no. 5) X/R is homeomorphic to the subspace \(s(\mathrm{X}/\mathrm{R})\) of X, which is Hausdorff. Furthermore \(s(\mathrm{X}/\mathrm{R})\) is the set of all \(x \in \mathbf{X}\) such that \(s(\varphi(x)) = x\) , where \(\varphi: \mathbf{X} \to \mathbf{X}/\mathbf{R}\) is the canonical mapping; hence the second assertion follows from no. 1, Proposition 2.

